\section{orders of infinity, asymptotic equivalence}
%%%%%
%%%%%

\begin{definition}[order of infinity/infinitesimal]%
  \label{def:ordine_infinito}%
  \mymark{***}%
  \mymargin{order of infinity/infinitesimal}%
\index{order of infinity/infinitesimal}%
  \index{order!of infinity}%
  \index{infinity}%
  \index{infinitesimal}%
  Let $A\subset \RR$, $f,g\colon A \to (0,+\infty)$.
  Let $x_0\in \bar \RR$ be an accumulation point of $A$.
  \begin{enumerate}
  \item
  We will say that
  for $x\to x_0$ the function $f$ is \emph{much smaller}
  than the function $g$ and we will write $f(x) \ll g(x)$ if
  \mymargin{$\ll$}%
\index{$\ll$}
  \[
  \frac{f(x)}{g(x)} \to 0, \qquad \text{as $x\to x_0$}
  \]
  we will instead say that $f$ is \emph{much larger}
  than $g$ and we will write $f(x) \gg g(x)$ if
  \mymargin{$\gg$}%
\index{$\gg$}
  \[
  \frac{f(x)}{g(x)} \to +\infty, \qquad \text{as $x\to x_0$.}
  \]
  \item
  Finally, we will say that $f$ and $g$
  are \emph{asymptotically equivalent}%
\mymargin{asymptotic equivalence}%
\index{asymptotically equivalent}
  \mymargin{asymptotic equivalence}%
\index{equivalence!asymptotic}%
  for $x\to x_0$
  and we will write $f(x) \sim g(x)$ if
  \mymargin{$\sim$}%
\index{$\sim$}
  \[
  \frac{f(x)}{g(x)} \to 1, \qquad \text{as $x\to x_0.$}
  \]
  \end{enumerate}
\end{definition}
  
For example, it is easy to verify that if $\alpha > \beta > 0$
then $x^\alpha \gg x^\beta$ for $x\to +\infty$
while $x^\alpha \ll x^\beta$ for $x\to 0^+$.
Similarly, if $a>b>1$ then $a^x\gg b^x$ for $x\to +\infty$
while $a^x \ll b^x$ for $x\to -\infty$.

It is very easy to verify that the relations
$\ll$ and $\gg$ are inverses of each other
and satisfy the transitive property,
while the relation $\sim$ satisfies the symmetric
and transitive properties.

\begin{theorem}[orders of infinity]
\label{th:ordine_infinito}%
\mymargin{orders of infinity}%
\index{orders of infinity}%
\mymark{***}%
Let $a>1$ and $\alpha>0$. For $n\to +\infty$, $n\in\NN$
we have
\[
  n^\alpha \ll a^n \ll n! \ll n^n
\]
and for $x\to +\infty$, $x\in \RR$ we have 
\[
\log_a x \ll x^\alpha \ll a^x.
\]
\end{theorem}
%
\begin{proof}
\mymark{**}
We begin by showing that $a^n \ll n!$
by applying the ratio test to the sequence $\frac{a^n}{n!}$:
\[
\frac{\displaystyle \frac{a^{n+1}}{(n+1)!}}{\displaystyle \frac{a^n}{n!}}
= \frac{a^{n+1}}{a^n}\cdot \frac{n!}{(n+1)!}
= a \cdot \frac {1}{n + 1} \to 0 < 1.
\]
Thus, as required, $a^n / n! \to 0$.
We proceed similarly to show that $n! \ll n^n$:
\begin{align*}
\frac{(n+1)!}{n!}\cdot \frac{n^n}{(n+1)^{n+1}}
&= (n+1) \cdot \enclose{\frac{n}{n+1}}^n \frac {1}{n+1}\\
&= \frac{1}{\enclose{1+\frac 1 n}^n} \to \frac 1 e < 1.
\end{align*}
  
To demonstrate that
$n^\alpha \ll a^n$
we can proceed with the ratio test, as in the previous cases:
\[
\frac{(n+1)^\alpha}{n^\alpha}\cdot \frac{a^n}{a^{n+1}}
= \frac 1 a \cdot \enclose{\frac{n+1}{n}}^\alpha \to \frac 1 a \cdot 1^\alpha =
\frac 1 a < 1
\]
from which $n^\alpha / a^n \to 0$.

For $x\in \RR$,
we try to reduce to a sequence with integer values.
We observe that
\[
\lfloor x \rfloor
\le x
\le \lfloor x \rfloor + 1
\]
from which, by monotonicity,
\[
\lfloor x \rfloor^\alpha
\le x^\alpha
\le (\lfloor x \rfloor + 1)^\alpha
= \lfloor x \rfloor^\alpha \enclose{1+ \frac{1}{\lfloor x \rfloor}}^\alpha
\]
and
\[
a^{\lfloor x \rfloor}
\le a^{x}
\le a^{\lfloor x \rfloor + 1}
= a \cdot a^{\lfloor x \rfloor}.
\]
Thus
\[
\frac{\lfloor x \rfloor^\alpha}{a \cdot a^{\lfloor x \rfloor}}
\le \frac{x^\alpha}{a^{x}}
\le \frac{\lfloor x \rfloor^\alpha \enclose{1+ \frac{1}{\lfloor x \rfloor}}^\alpha}
    {a^{\lfloor x \rfloor}}.
\]
But now, if $x\to +\infty$ knowing that $n = \lfloor x\rfloor \to +\infty$ 
we can perform a change of variable in the limit
\[
\lim_{x\to +\infty} \frac{\lfloor x \rfloor^\alpha}{a^{\lfloor x \rfloor}} 
= \lim_{n\to+\infty} \frac{n^\alpha}{a^n} = 0
\]
from which it follows that $\frac{x^\alpha}{a^{x}}\to 0$.

To demonstrate the last relation, $\log_a x\ll x^\alpha$,
we perform the change of variable $y = \alpha \cdot \log_a x$
so that $a^y = x^\alpha$.
Note that if $x\to +\infty$
then $y \to +\infty$.
Thus, by the previous properties,
we know that $y \ll a^y$ and therefore
\[
\frac{\log_a x}{x^\alpha}
= \frac{1}{\alpha}\cdot\frac{y}{a^{y}} \to 0.
\]
\end{proof}

The notations and
orders of infinity identified in the previous theorem
are very useful tools in the calculation of limits.

Asymptotic equivalence
is preserved under multiplication and division:
if $f\sim F$ and $g\sim G$ then
\[
 f \cdot g \sim F \cdot G,
 \qquad
 \frac{f}{g} \sim \frac{F}{G}.
\]
We also observe that if
$f \sim g$ and if $f\to \ell$ then
also $g\to \ell$.
If $\ell\in(0,+\infty)$
the relation $f\sim \ell$ is equivalent to $f\to \ell$.

Regarding addition,
it is easy to verify that if $f\ll g$ then
$(f+g) \sim g$ since
\[
  \frac{f + g}{g} = \frac{f}{g} + 1 \to 1.
\]

In a limit involving sums of terms
of different orders, we can then factor out the terms of maximum order to
identify the limit, as we do in the following.

\begin{example}
Calculate the limit
\[
\lim_{n\to+\infty}
\frac{2n^4 + 3^n - 3 \ln n}{n! - 3\sqrt n}.
\]
\end{example}
\begin{proof}[Solution.]
We have
\[
\frac{2n^4 + 3^n - 3 \ln n}{n! - 3\sqrt{n}}
= \frac
{3^n \cdot \enclose{2\frac{n^4}{3^n}+ 1 - 3\frac{\ln n}{3^n}}}
{n!\cdot \enclose{1-3\frac{\sqrt n}{n!}}}
\]
and recalling that (theorem~\ref{th:ordine_infinito})
\[
n^4 \ll 3^n, \qquad
\ln n \ll 3^n, \qquad
\sqrt n \ll n!, \qquad
3^n \ll n!
\]
we will have
\[
\frac{2n^4 + 3^n - 3 \ln n}{n! - 3\sqrt{n}}
\sim \frac{3^n}{n!} \to 0.
\]
\end{proof}


\begin{exercise}
Calculate the following limits
\begin{gather*}
  \lim_{n\to +\infty} \frac{\displaystyle \ln\sqrt{n^2+n^n}}
  {\displaystyle e^{1 + \ln n}\cdot \ln(n^2-n\sqrt n)}, \qquad
  \lim_{n\to +\infty} \frac{\sqrt{n! + 2^n}}{3^n}, \\
  \lim_{n\to +\infty} \frac{\sqrt{(2n)!}}{n^n}, \qquad
  \lim_{n\to +\infty} \sqrt[n]{e^n + \sqrt{10^n}}.
\end{gather*}
\end{exercise}

\begin{exercise}
  Prove that $\sqrt[n]{n}\to 1$ as $n\to +\infty$.
\end{exercise}



%%%%%%%%%%%
%%%%%%%%%%%